\SourceLeaf{330}%
\[
\begin{aligned}
g&=-\frac{k^3\cos T}{a(1-ee)}\\
g'&=+\frac{k^3\sin T}{a(1-ee)}\\
g''&=+\frac{3kee\cos T}{a(1-ee)}\mathord{.}cs\\
g'''&=+\frac{k^3\cos T}{a(1-ee)}\\
g''''&=-\frac{3kee\sin T}{a(1-ee)}\mathord{.}cs\\
g^{\mathrm{v}}&=+\frac{3ee\cos T}{a(1-ee)k}\mathord{.}(1-(2+2ee)ss+3ees^4)\\
h&=-\frac{k\sin T}{ac}\\
h'&=-\frac{k\cos T}{ac}\\
h''&=-\frac{\sin T}{acck}\mathord{.}(1-ee)s\\
h'''&=+\frac{k\sin T}{ac}\\
h''''&=-\frac{\cos T}{acck}\mathord{.}(1-ee)s\\
h^{\mathrm{v}}&=-\frac{(1-ee)\sin T}{ac^3k^3}\mathord{.}(1+ss-2ees^4)
\end{aligned}
\]

\Article{29}
Wir wollen nun die drei Gleichungen (7), (8), (6) in folgende Form setzen
\[
\begin{aligned}
t&=-fr(1+Frr)\\
b&=-gr(1+Grr)\\
l&=-hr(1+Hrr)
\end{aligned}
\]
wo
\[
\begin{aligned}
-fr&=\frac{k\sin T\mathord{.}\operatorname{tang}B}{a}\mathord{.}r\\
-gr&=\frac{k^3\cos T}{a(1-ee)}\mathord{.}r\\
-hr&=\frac{k\sin T}{a\cos B}\mathord{.}r
\end{aligned}
\]
beziehungsweise die genäherten und bis auf die dritte Ordnung ausschl. genauen Werthe von \(t,b,l\) sind, die zur Abkürzung mit \(\tau,\beta,\lambda\) bezeichnet werden sollen. Jede der Grössen \(F,G,H\) ist das Aggregat von sieben Theilen, nemlich
